
\section{\BookRecordIndex{T0142}格罗莫夫拓扑}\label{sec:5-6}
\zhindex{geluomofutuopu}{格罗莫夫拓扑}

我们希望利用定义~\ref{def:5-5-1} 在 $\overline{\mathcal M}_{0,n}(M,A;J)$ 上定义一个
称为\emph{\BookRecordIndex{T0142}格罗莫夫拓扑}的拓扑。在这个拓扑下，等价类序列
$[\mathbf u^\nu,\mathbf z^\nu]$ 应当收敛到 $[\mathbf u,\mathbf z]$，只要
$(\mathbf u^\nu,\mathbf z^\nu)$ \BookRecordIndex{T0141}格罗莫夫收敛到 $(\mathbf u,\mathbf z)$。若子集
$F\subset\overline{\mathcal M}_{0,n}(M,A;J)$ 中每个\BookRecordIndex{T0141}格罗莫夫收敛列的极限仍属于 $F$，
就称 $F$ \emph{\BookRecordIndex{T0138}格罗莫夫闭}。若子集
$U\subset\overline{\mathcal M}_{0,n}(M,A;J)$ 的补集\BookRecordIndex{T0138}格罗莫夫闭，就称 $U$
\emph{\BookRecordIndex{T0140}格罗莫夫开}。这些集合显然在 $\overline{\mathcal M}_{0,n}(M,A;J)$ 上定义一个
拓扑。不过，要证明关于这个拓扑的收敛等价于\BookRecordIndex{T0141}格罗莫夫收敛，还需要一些准备。
\zhindex{geluomofubi}{格罗莫夫闭}
\zhindex{geluomofukai}{格罗莫夫开}

\paragraph{拓扑与序列。}
设 $(X,\mathcal U)$ 是\BookRecordIndex{T0102}第一可数的\BookRecordIndex{T0163}豪斯多夫拓扑空间（也就是说，$X$ 中每一点都有可数邻域
基）。在这些假设下，极限是唯一的；而且，拓扑 $\mathcal U$ 完全由收敛序列描述。换言之，
子集 $A\subset X$ 的闭包 $\operatorname{cl}(A)$ 恰好由 $A$ 中收敛序列的所有极限组成。
详细说来，记
\[
  C=C(\mathcal U)\subset X\times X^{\mathbb N}
\]
为所有二元组 $(x_0,(x_n)_n)$ 所成的集合，其中 $x_0\in X$，$(x_n)_n$ 是 $X$ 中收敛到
$x_0$ 的序列。收敛序列族 $C$ 具有下列性质。

\emph{（常值）}若对所有 $n\in\mathbb N$ 都有 $x_n=x_0$，则
$(x_0,(x_n)_n)\in C$。

\emph{（子列）}若 $(x_0,(x_n)_n)\in C$，且 $g:\mathbb N\to\mathbb N$ 严格递增，则
$(x_0,(x_{g(n)})_n)\in C$。

\emph{（子列的子列）}若对每个严格递增函数 $g:\mathbb N\to\mathbb N$，都存在严格递增
函数 $f:\mathbb N\to\mathbb N$，使 $(x_0,(x_{g\circ f(n)})_n)\in C$，则
$(x_0,(x_n)_n)\in C$。

\emph{（对角）}若 $(x_0,(x_k)_k)\in C$，且对每个 $k$ 都有
$(x_k,(x_{k,n})_n)\in C$，则存在序列 $k_i,n_i\in\mathbb N$，使
$(x_0,(x_{k_i,n_i})_i)\in C$。

\emph{（唯一性）}若 $(x_0,(x_n)_n)\in C$ 且 $(y_0,(x_n)_n)\in C$，则 $x_0=y_0$。

现在反过来，从任意序列族 $C\subset X\times X^{\mathbb N}$ 出发，尝试在 $X$ 上定义
拓扑。给定 $C$，定义
\[
  \mathcal U(C)\subset 2^X
\]
为所有满足下述条件的子集 $U\subset X$ 所成的集合：
\[
  (x_0,(x_n)_n)\in C\cap(U\times X^{\mathbb N})
  \quad\Longrightarrow\quad
  \exists N\in\mathbb N\ \forall n\in\mathbb N
  \bigl(n\ge N\Longrightarrow x_n\in U\bigr).
\]
下面的练习~\ref{ex:5-6-1} 说明 $\mathcal U(C)$ 总是一个拓扑。不过，我们希望
$\mathcal U(C)$ 还具有如下性质：关于 $\mathcal U(C)$ 收敛的序列集合
$C(\mathcal U(C))$ 恰好就是原来的集合 $C$。我们将看到，（子列）公理保证对每个子集
$F\subset X$ 都有
\[
  X\setminus F\in\mathcal U(C)
  \quad\Longleftrightarrow\quad
  \bigl((x_0,(x_n)_n)\in C\cap(X\times F^{\mathbb N})
  \Longrightarrow x_0\in F\bigr).
\]
为了证明子集 $A\subset X$ 的闭包由所有满足如下条件的 $x_0\in X$ 组成：存在序列
$x_n\in A$ 使 $(x_0,(x_n)_n)\in C$，还需要（常值）与（对角）公理。要证明
$C(\mathcal U(C))=C$，则五条公理全都需要。

\begin{exercise}\label{ex:5-6-1}
设 $C$ 是 $X\times X^{\mathbb N}$ 的任意子集。证明 $\mathcal U(C)$ 是一个拓扑。证明
\[
  C\subset C(\mathcal U(C)),
  \qquad
  \mathcal U(C)=\mathcal U(C(\mathcal U(C))).
\]
\end{exercise}

\begin{exercise}\label{ex:5-6-2}
设 $C$ 满足（子列）公理，并令 $\mathcal F\subset 2^X$ 是所有满足
\[
  (x_0,(x_n)_n)\in C\cap(X\times F^{\mathbb N})
  \quad\Longrightarrow\quad x_0\in F
\]
的子集 $F\subset X$ 所成的集合。证明
\[
  F\in\mathcal F
  \quad\Longleftrightarrow\quad
  X\setminus F\in\mathcal U(C).
\]
\emph{提示：}只有在证明 $F\in\mathcal F\Longrightarrow X\setminus F\in\mathcal U(C)$
时才需要（子列）公理。
\end{exercise}

\begin{lemma}\label{lem:5-6-3}
设 $C$ 满足（常值）、（子列）和（对角）公理，$A\subset X$ 是任意子集。则 $A$ 关于
$\mathcal U(C)$ 的闭包为
\[
  \operatorname{cl}(A)
  =\bigl\{x_0\in X\bigm|
    \exists(x_n)_n\in A^{\mathbb N}\text{ 使 }(x_0,(x_n)_n)\in C\bigr\}.
\]
\end{lemma}

\begin{proof}
令 $F$ 为等号右端的集合。（常值）公理表明 $A\subset F$，而且由定义可知，每个包含 $A$
的闭集都包含 $F$。由练习~\ref{ex:5-6-2} 与（对角）公理，$F$ 是闭集。这就证明了引理。
\end{proof}

\begin{lemma}\label{lem:5-6-4}
设 $C$ 是 $X$ 中满足（常值）、（子列）、（子列的子列）、（对角）与（唯一性）公理的
序列族。那么，$X$ 上存在唯一拓扑 $\mathcal U$
\sourcefootnote{原书断言拓扑“唯一”，但随后的证明只构造并验证了序列化拓扑
$\mathcal U(C)$，即证明 $C(\mathcal U(C))=C$；一般拓扑未必由其收敛序列唯一决定。若把
范围限定为由序列刻画的序列型拓扑，则可得到相应的唯一性，但原文未加这一限制。}，使 $C$
恰为关于它收敛的序列集合。
\end{lemma}

\begin{proof}
由前面的讨论，只须验证 $C=C(\mathcal U(C))$。练习~\ref{ex:5-6-1} 给出
$C\subset C(\mathcal U(C))$，所以还须证明 $C(\mathcal U(C))\subset C$。用反证法，假设
$(x_0,(x_n)_n)\notin C$。由（子列的子列）公理，存在严格递增函数
$g:\mathbb N\to\mathbb N$，使得对每个严格递增函数 $f:\mathbb N\to\mathbb N$ 都有
\begin{equation}\label{eq:5-6-1}
  (x_0,(x_{g\circ f(n)})_n)\notin C.
\end{equation}
由（常值）公理，这蕴含存在 $N\in\mathbb N$，使对所有 $n\in\mathbb N$ 都有
\begin{equation}\label{eq:5-6-2}
  n\ge N\quad\Longrightarrow\quad x_{g(n)}\ne x_0.
\end{equation}
定义
\[
  F:=\operatorname{cl}\bigl(\{x_{g(n)}\mid n\ge N\}\bigr).
\]
我们断言 $x_0\notin F$。反之，若 $x_0\in F$，则由引理~\ref{lem:5-6-3}，存在整数序列
$n_i\ge N$，使 $(x_0,(x_{g(n_i)})_i)\in C$。若 $(n_i)_i$ 有界，由（子列）公理，不妨
假设对所有 $i$ 都有 $n_i=n\ge N$。由（常值）公理，
$(x_{g(n)},(x_{g(n_i)})_i)\in C$；于是（唯一性）公理给出
$x_0=x_{g(n)}$，与式~\eqref{eq:5-6-2} 矛盾。因此 $(n_i)_i$ 无界。再由（子列）公理，不妨
假设对所有 $i$ 都有 $n_{i+1}>n_i$。这样就找到了严格递增函数
$f:\mathbb N\to\mathbb N$，使 $(x_0,(x_{g\circ f(i)})_i)\in C$，与
式~\eqref{eq:5-6-1} 矛盾。因此 $x_0\notin F$，断言得证。于是我们找到了 $x_0$ 的开邻域
$U:=X\setminus F$ 与严格递增函数 $g:\mathbb N\to\mathbb N$，使得对所有 $n\ge N$ 都有
$x_{g(n)}\notin U$。所以 $(x_n)_n$ 关于拓扑 $\mathcal U(C)$ 不收敛到 $x_0$。引理得证。
\end{proof}

为了更好地理解拓扑 $\mathcal U(C)$，例如判断它何时是\BookRecordIndex{T0163}豪斯多夫的，我们需要知道它\BookRecordIndex{T0102}第一
可数，即 $X$ 中每一点都有可数邻域基。这还需要下面引理~\ref{lem:5-6-5} 中的进一步假设。

\paragraph{\BookRecordIndex{T0142}格罗莫夫拓扑。}
考虑\BookRecordIndex{T0340}稳定映射等价类空间 $X=\overline{\mathcal M}_{0,n}(A;J)$。\BookRecordIndex{T0340}稳定映射列的\BookRecordIndex{T0141}格罗莫夫收敛
定义具有这样的形式：它能下降到 $\overline{\mathcal M}_{0,n}(A;J)$。换言之，给定序列
$\mathbf x^\nu$ 与极限候选 $\mathbf x$，序列 $\mathbf x^\nu$ \BookRecordIndex{T0141}格罗莫夫收敛到
$\mathbf x$，当且仅当任取分别与 $\mathbf x^\nu$ 等价的\BookRecordIndex{T0340}稳定映射 $\mathbf y^\nu$ 以及
与 $\mathbf x$ 等价的\BookRecordIndex{T0340}稳定映射 $\mathbf y$，序列 $\mathbf y^\nu$ 都\BookRecordIndex{T0141}格罗莫夫收敛到
$\mathbf y$。因此，
讨论 $\overline{\mathcal M}_{0,n}(A;J)$ 中的\BookRecordIndex{T0141}格罗莫夫收敛列
$[\mathbf x^\nu]$ 是有意义的。此外，由注~\ref{rem:5-5-7} 与引理~\ref{lem:5-5-8}、
引理~\ref{lem:5-5-9}，这种收敛可以用一族函数 $\rho_\varepsilon(\mathbf x,\mathbf x')$ 来
描述，而这些函数满足下一个引理中的性质。因此，可以把下一个引理的结论用于
$X=\overline{\mathcal M}_{0,n}(A;J)$，并令 $C$ 为所有\BookRecordIndex{T0141}格罗莫夫收敛列
$([\mathbf x],[\mathbf x^\nu]_\nu)$ 所成的集合。

\begin{lemma}\label{lem:5-6-5}
设 $X$ 是集合，$C\subset X\times X^{\mathbb N}$ 是 $X$ 中满足（唯一性）公理的序列族。
假设对每个 $x\in X$，都存在常数 $\varepsilon_0(x)>0$，而且对每个
$0<\varepsilon<\varepsilon_0(x)$ 都有一族函数
\[
  X\longrightarrow[0,\infty],
  \qquad x'\longmapsto\rho_\varepsilon(x,x'),
\]
满足下列条件。
\begin{enumerate}[label=(\alph*),leftmargin=3em]
\item 若 $x\in X$ 且 $0<\varepsilon<\varepsilon_0(x)$，则
  $\rho_\varepsilon(x,x)=0$。
\item 若 $x\in X$、$0<\varepsilon<\varepsilon_0(x)$ 且
  $(x_n)_n\in X^{\mathbb N}$，则
  \[
    (x,(x_n)_n)\in C
    \quad\Longleftrightarrow\quad
    \lim_{n\to\infty}\rho_\varepsilon(x,x_n)=0.
  \]
\item 若 $x\in X$、$0<\varepsilon<\varepsilon_0(x)$ 且
  $(x',(x_n)_n)\in C$，则
  \[
    \rho_\varepsilon(x,x')<\varepsilon
    \quad\Longrightarrow\quad
    \limsup_{n\to\infty}\rho_\varepsilon(x,x_n)
    \le\rho_\varepsilon(x,x').
  \]
\end{enumerate}
那么 $C=C(\mathcal U(C))$。此外，拓扑 $\mathcal U(C)$ \BookRecordIndex{T0102}第一可数且为\BookRecordIndex{T0163}豪斯多夫拓扑。
\end{lemma}

\begin{proof}
由 (a) 与 (b)，$C$ 满足（常值）、（子列）和（子列的子列）公理。为证明（对角）公理，
假设 $(x,(x_k)_k)\in C$，而且对 $k\ge1$ 有 $(x_k,(x_{k,n})_n)\in C$。选取实数
$0<\varepsilon<\varepsilon_0(x)$。由 (b)，存在整数序列 $(k_i)_i$，使
\[
  \rho_\varepsilon(x,x_{k_i})\le2^{-i}\varepsilon.
\]
再由 (c)，存在整数序列 $(n_i)_i$，使
\[
  \rho_\varepsilon(x,x_{k_i,n_i})\le2^{1-i}\varepsilon.
\]
所以由 (b)，$(x,(x_{k_i,n_i})_i)\in C$。这说明 $C$ 满足全部五条公理；由
引理~\ref{lem:5-6-4}，$C=C(\mathcal U(C))$。

下面证明 $X$ 中每一点都有可数邻域基。给定 $x$，选取实数
$0<\varepsilon<\varepsilon_0(x)$，并考虑
\[
  U_n(x):=\{x'\in X\mid\rho_\varepsilon(x,x')<\varepsilon/n\},
  \qquad n\in\mathbb N.
\]
我们断言这些集合构成 $x$ 的开邻域基。先注意，若 $y_0\in U_n(x)$ 且
$(y_0,(y_k)_k)\in C$，则由 (c)
\[
  \limsup_{k\to\infty}\rho_\varepsilon(x,y_k)
  \le\rho_\varepsilon(x,y_0)<\varepsilon/n.
\]
所以充分大的 $k$ 都满足 $y_k\in U_n(x)$。由 $\mathcal U(C)$ 的定义，$U_n(x)$ 是开集。
特别地，对每个 $x\in X$ 与每个 $n\in\mathbb N$，$U_n(x)$ 都是 $x$ 的邻域。为证明
$U_n(x)$ 构成邻域基，用反证法。若不然，就存在 $x$ 的开邻域 $U$，它不包含任何
$U_n(x)$。于是对每个 $n$，都有一点 $x_n\in U_n(x)\setminus U$。由 (b)，
$(x,(x_n)_n)\in C$。因为 $U$ 是包含 $x$ 的开集，充分大的 $n$ 必有 $x_n\in U$，这与
假设矛盾。

由此已证明 $\mathcal U(C)$ \BookRecordIndex{T0102}第一可数。为了证明它是\BookRecordIndex{T0163}豪斯多夫的，须说明任意两个不同点
$x,y$ 都有不交的开邻域。我们证明：若 $x\ne y$，则存在整数
$n=n(x,y)$，使 $U_n(x)\cap U_n(y)=\varnothing$。否则就存在序列 $(x_n)_n\subset X$，使
对所有 $n$ 都有 $x_n\in U_n(x)\cap U_n(y)$。于是由 (b)，序列 $(x_n)_n$ 同时收敛到
$x$ 和 $y$，这与极限唯一的假设矛盾。引理得证。
\end{proof}

现在把\BookRecordIndex{T0246}模空间 $\overline{\mathcal M}_{0,n}(M,A;J)$ 上的\BookRecordIndex{T0142}格罗莫夫拓扑定义为
$\mathcal U(C)$，其中 $C$ 是所有\BookRecordIndex{T0141}格罗莫夫收敛列所成的集合。下一个定理说明，这个拓扑
具有所有预期性质。

\begin{theorem}\label{thm:5-6-6}
\begin{enumerate}[label=(\roman*),leftmargin=3em]
\item $\overline{\mathcal M}_{0,n}(M,A;J)$ 中的序列关于\BookRecordIndex{T0142}格罗莫夫拓扑收敛，当且仅当它
  \BookRecordIndex{T0141}格罗莫夫收敛。
\item $\overline{\mathcal M}_{0,n}(M,A;J)$ 是紧且\BookRecordIndex{T0212}可度量化的空间。特别地，它是\BookRecordIndex{T0163}豪斯多夫
  空间，并有可数稠密子集。
\item \BookRecordIndex{T0283}求值映射
  \[
    \operatorname{ev}_i:\overline{\mathcal M}_{0,n}(M,A;J)\longrightarrow M,
  \]
  投影
  \[
    \pi:\overline{\mathcal M}_{0,n}(M,A;J)
    \longrightarrow\overline{\mathcal M}_{0,n-1}(M,A;J),
  \]
  以及第~\ref{sec:5-1} 节的\BookRecordIndex{T0377}遗忘映射
  \[
    \pi:\overline{\mathcal M}_{0,n}(M,A;J)
    \longrightarrow\overline{\mathcal M}_{0,n}
  \]
  关于\BookRecordIndex{T0142}格罗莫夫拓扑连续。
\end{enumerate}
\end{theorem}

\begin{proof}
由定理~\ref{thm:5-5-3}，\BookRecordIndex{T0141}格罗莫夫收敛列的极限唯一。因此，由注~\ref{rem:5-5-7} 与
引理~\ref{lem:5-5-8}、\ref{lem:5-5-9}，\BookRecordIndex{T0141}格罗莫夫收敛列满足引理~\ref{lem:5-6-5} 的全部
假设。这证明了 (i)，并说明\BookRecordIndex{T0142}格罗莫夫拓扑\BookRecordIndex{T0102}第一可数且为\BookRecordIndex{T0163}豪斯多夫拓扑。

我们证明 $\overline{\mathcal M}_{0,n}(M,A;J)$ 上的\BookRecordIndex{T0142}格罗莫夫拓扑有可数基。为此，固定
$n$-\BookRecordIndex{T0084}带标号树 $T=(T,E,\Lambda)$。所有以 $T$ 为模型的
$\mathbf x\in\mathcal{SC}_{0,n}(M,A;J)$ 所成的集合
$\widetilde{\mathcal M}_{0,T}(M,A;J)$，是某个\BookRecordIndex{T0213}可分 Banach 流形的可数多个紧子集之并，
因而含有可数稠密子集。所以，每个空间 $\widetilde{\mathcal M}_{0,T}(M,A;J)$ 都有可数
稠密子集，其商\BookRecordIndex{T0072}层
\[
  \mathcal M_{0,T}(M,A;J)
  =\widetilde{\mathcal M}_{0,T}(M,A;J)/G_T
\]
也有可数稠密子集。由于\BookRecordIndex{T0072}层只有有限多个，
$\overline{\mathcal M}_{0,n}(M,A;J)$ 也有可数稠密子集。因此，
$\overline{\mathcal M}_{0,n}(M,A;J)$ 上的\BookRecordIndex{T0142}格罗莫夫拓扑有可数稠密子集。把这个可数稠密
子集内各点的可数邻域基取并，便得到\BookRecordIndex{T0142}格罗莫夫拓扑的可数基
\sourcefootnote{原书此处由\BookRecordIndex{T0102}第一可数性与可数稠密子集直接\BookRecordIndex{T0326}推出\BookRecordIndex{T0101}第二可数性；一般而言，这一
推理并不成立。要得到所声称的可数基，还需另行证明该\BookRecordIndex{T0142}格罗莫夫拓扑\BookRecordIndex{T0101}第二可数。}。

其次，由定理~\ref{thm:5-5-5}，$\overline{\mathcal M}_{0,n}(M,A;J)$ 中每个序列都有
\BookRecordIndex{T0141}格罗莫夫收敛子列。因此，\BookRecordIndex{T0142}格罗莫夫拓扑\BookRecordIndex{T0371}序列紧。每个有可数基的\BookRecordIndex{T0371}序列紧拓扑空间都是紧空间
（\srcite{338}，第 9 章第 2 节），故
$\overline{\mathcal M}_{0,n}(M,A;J)$ 是紧空间。\footnote{这里给出一个直接论证。设 $X$
是有可数基的\BookRecordIndex{T0371}序列紧拓扑空间。于是每个开覆盖都有可数加细
$\{U_n\}_{n\in\mathbb N}$，它由基中那些包含于原覆盖某个成员的元素组成。假设这个加细
没有有限子覆盖。取一个仍覆盖 $X$ 的子列后，不妨对每个 $n$ 都有
$U_n\not\subset U_1\cup\cdots\cup U_{n-1}$。因而可取序列 $x_n\in U_n$，使得对 $m<n$
都有 $x_n\notin U_m$。取收敛子列 $x_{n_i}\to x$。于是对某个 $m$ 有 $x\in U_m$，从而
存在 $i_0\in\mathbb N$，使得 $i\ge i_0$ 时 $x_{n_i}\in U_m$。取一个满足 $n_i>m$ 的
$i$，便得到矛盾。因此，该加细有有限子覆盖，原开覆盖也有有限子覆盖。}
现在，每个紧\BookRecordIndex{T0163}豪斯多夫空间都是正规空间（不交闭子集可用开集分离）。Urysohn \BookRecordIndex{T0212}可度量化定理
\srcite{338}（第 8 章第 3 节）断言，每个有可数基的正规拓扑空间都\BookRecordIndex{T0212}可度量化。这证明了
(ii)。

\BookRecordIndex{T0283}求值映射的连续性直接由定义得出。\BookRecordIndex{T0377}遗忘映射
$\overline{\mathcal M}_{0,n}(M,A;J)\to\overline{\mathcal M}_{0,n}$ 的连续性由
定理~\ref{thm:5-2-2}\textup{(iv)} 与命题~\ref{prop:5-5-2} 得出。投影
$\overline{\mathcal M}_{0,n}(M,A;J)\to
\overline{\mathcal M}_{0,n-1}(M,A;J)$ 的连续性证明留作练习。定理得证。
\end{proof}
